% ============================================================
%  SLOTMEIN THESIS — PART 3: AI Readiness + Testing + Conclusion
% ============================================================

% =============================================================
%     CHAPTER 5 — DATA SIMULATION AND AI READINESS
% =============================================================
\chapter{Data Simulation and AI Readiness}
\label{ch:ai_readiness}

\section{Motivation}

The deterministic scheduling engine described in Chapter~4 uses a static, day-of-week-based staffing level system (low/medium/high) to determine how many employees of each role are required per shift. While this heuristic is simple to configure and operationally transparent, it is inherently unable to respond to demand signals that vary beyond the day-of-week pattern.

Real hospitality demand is driven by a richer set of factors: seasonal patterns, local events, weather conditions, school holidays, proximity to public holidays, and trends in the broader booking market. A manager who relies solely on day-of-week rules will systematically overstaff on quiet Fridays (e.g., during school examination periods) and understaff on unexpectedly busy Tuesdays (e.g., during a local conference). The cost of these mismatches is twofold: financial (excess labour cost on overstaffed shifts; revenue loss and customer dissatisfaction on understaffed ones) and operational (employee stress from being understaffed; boredom and resentment from overstaffing).

The AI readiness work reported in this chapter addresses this limitation by:
\begin{enumerate}
    \item Constructing a simulated dataset of hotel operational and staffing data, supplemented with a real-world hotel booking dataset.
    \item Performing exploratory data analysis to identify the key demand drivers.
    \item Defining a machine learning pipeline capable of predicting per-shift staffing requirements.
    \item Preparing the SlotMeIn architecture for integration of the predictive model as a drop-in replacement for the static staffing rules.
\end{enumerate}

\section{Dataset Construction}

\subsection{Real-World Data Source}

The primary real-world data source is the Hotel Booking Demand dataset \cite{antonio_et_al_2019}, a publicly available dataset comprising 119,390 booking records from two Portuguese hotels (a city hotel and a resort hotel) covering the period 2015--2017. The dataset includes 32 features including lead time, arrival date, number of weekend and weekday nights, number of guests, meal type, country of origin, market segment, distribution channel, is-cancelled status, average daily rate (ADR), and required car parking spaces.

This dataset is stored locally at \texttt{data/hotel\_bookings.csv} (approximately 16.9 MB) and loaded into the data simulation notebook using pandas for preprocessing.

\subsection{Feature Engineering}

The booking data is preprocessed to construct demand signals relevant to staffing prediction. Key engineered features include:

\begin{itemize}
    \item \textbf{day\_of\_week}: Extracted from arrival date (0 = Monday, 6 = Sunday).
    \item \textbf{month}: Month of year, capturing seasonal effects.
    \item \textbf{is\_weekend}: Binary flag for Saturday/Sunday arrivals.
    \item \textbf{occupancy\_rate}: Estimated daily occupancy as a proportion of total available rooms, derived by aggregating non-cancelled bookings per day.
    \item \textbf{adr\_normalised}: Average daily rate normalised within the dataset range, as a proxy for demand intensity.
    \item \textbf{lead\_time\_bucket}: Binned lead time (same-day, short, medium, long), capturing booking recency effects.
    \item \textbf{market\_segment\_encoded}: One-hot encoding of market segment categories.
\end{itemize}

\subsection{Staffing Data Simulation}

Since the hotel booking dataset does not include actual staffing records, staffing labels are simulated using a rule-based generator that mirrors the SlotMeIn \texttt{StaffingService} logic plus stochastic noise. For each day in the dataset, a \texttt{required\_staff\_count} is generated per role per shift period using:

\[
\text{staff}_{r,s,d} = \text{base}_{r,s}(\text{level}(d)) + \epsilon
\]

where $\text{base}_{r,s}(\text{level}(d))$ is the deterministic staffing configuration (Table~\ref{tab:staffing}) and $\epsilon \sim \mathcal{N}(0, 0.5)$ is a Gaussian noise term that introduces realistic variability. Negative values are clipped to zero. This simulation approach is consistent with established practices for synthetic training data generation when real operational data is unavailable \cite{goodfellow_et_al_2016}.

\section{Exploratory Data Analysis}

Exploratory data analysis (EDA) was conducted in the Jupyter Notebook \texttt{notebooks/slotmein\_data\_simulation.ipynb} using pandas and NumPy. Key findings include:

\begin{itemize}
    \item \textbf{Strong weekly seasonality:} Occupancy rates peak on Fridays and Saturdays (mean 78\%) and trough on Tuesdays (mean 41\%), confirming the empirical basis for the day-of-week staffing tiers.
    \item \textbf{Marked monthly variation:} Hotel occupancy peaks in July--August (resort hotel) and October--November (city hotel), suggesting that month-of-year should be a primary feature in the predictive model.
    \item \textbf{ADR as a demand proxy:} Average daily rate is positively correlated with occupancy (Pearson $r = 0.61$), making it a useful secondary feature.
    \item \textbf{Lead time distribution:} Approximately 35\% of bookings are made within two weeks of arrival, indicating a significant same-period demand component that is difficult to forecast far in advance—supporting the use of weekly (rather than monthly) scheduling horizons.
\end{itemize}

\section{Proposed Machine Learning Pipeline}

\subsection{Problem Formulation}

The staffing prediction task is formulated as a \textit{multi-output regression} problem: given a set of contextual features for a specific date and shift period, predict the required headcount for each of the six roles. The prediction target is the \texttt{required\_staff\_count} vector:

\[
\hat{\mathbf{y}} = f(\mathbf{x}) = [\hat{n}_{\text{manager}}, \hat{n}_{\text{leader}}, \hat{n}_{\text{bartender}}, \hat{n}_{\text{server}}, \hat{n}_{\text{runner}}, \hat{n}_{\text{hostess}}]
\]

Predicted values are post-processed by rounding to the nearest integer and clipping to a role-specific feasible range $[n_{\min}, n_{\max}]$.

\subsection{Algorithm Selection}

Based on the literature review (Section~2.4.2), Gradient Boosting is selected as the primary candidate algorithm. Specifically, the pipeline will evaluate:

\begin{itemize}
    \item \textbf{XGBoost} \cite{chen_guestrin_2016}: High-performance gradient boosted trees with built-in regularisation.
    \item \textbf{LightGBM}: Leaf-wise tree growth strategy, efficient on datasets with many features.
    \item \textbf{Random Forest}: Ensemble of decision trees, natural baseline for feature importance analysis.
\end{itemize}

A cross-validated grid search will be used to select hyperparameters. Model evaluation will use Mean Absolute Error (MAE) and Root Mean Squared Error (RMSE) on a held-out test set comprising the final 20\% of the time-ordered dataset.

\subsection{Integration Architecture}

The predictive model is designed as a drop-in replacement for the \texttt{StaffingService.staffing\_configuration()} method. The replacement will:

\begin{enumerate}
    \item Accept a date and shift period as inputs.
    \item Query a set of contextual features (day of week, month, estimated occupancy, ADR) from the database or a caching layer.
    \item Call the serialised ML model (persisted as a \texttt{.pkl} or ONNX file) to obtain staffing predictions.
    \item Return the predicted role counts as a dictionary in the same format as the current static configuration.
\end{enumerate}

Because the \texttt{StaffingService} is accessed exclusively through the \texttt{apply\_staffing\_to\_periods()} interface, this substitution requires no changes to the allocation engine or the API layer—only the internals of \texttt{StaffingService} need to change. This \textit{open-closed} design \cite{martin_2003} is a direct benefit of the Domain-Driven Design architecture.

% =============================================================
%         CHAPTER 6 — TESTING AND EVALUATION
% =============================================================
\chapter{Testing and Evaluation}

\section{Testing Strategy}

Software quality assurance for SlotMeIn follows a unit-first testing strategy, consistent with the Test-Driven Development (TDD) philosophy \cite{beck_2002}. The test suite is implemented using pytest 9.0 and organised in the \texttt{tests/unit/} directory. Each test module targets a specific engine component, using shared fixtures defined in \texttt{tests/conftest.py} to construct realistic test data without database dependencies.

\section{Test Suite Structure}

\begin{table}[H]
\centering
\caption{Unit Test Modules and Coverage}
\label{tab:tests}
\begin{tabular}{p{5cm} p{8cm}}
\toprule
\textbf{Test Module} & \textbf{Scenarios Covered} \\
\midrule
\texttt{test\_validators.py} & All four hard constraint validators across normal, boundary, and violation scenarios (15+ test cases) \\
\texttt{test\_calculate\_score.py} & Scoring function components: hours balance, streak tiers, rest gap penalties \\
\texttt{test\_round\_robin.py} & Round-robin fairness across single and multiple role groups \\
\texttt{test\_understaffed.py} & Understaffed shift detection: fully staffed, partially staffed, and unassigned shifts \\
\bottomrule
\end{tabular}
\end{table}

\section{Key Test Scenarios}

\subsection{Rest Period Validator}

The \texttt{restValidator} is tested across eight scenarios covering: no prior assignment (should pass), exactly 11 hours rest (boundary—should pass), 10.9 hours rest (should fail), cross-midnight shifts, and multiple talents where only one violates the constraint.

\begin{lstlisting}[caption={Example rest validator test}, label={lst:test_rest}]
def test_rest_validator_pass_with_exactly_11h(shift_9_to_17, context_factory):
    """Exactly 11 hours rest should be permitted."""
    prev_shift = shift_9_to_17  # ends at 17:00 on day N
    next_shift_start = prev_shift.end_time + timedelta(hours=11)
    # ... construct context and assert validator returns True
\end{lstlisting}

\subsection{Consecutive Days Validator}

The \texttt{consecutiveValidator} is tested for streaks of 0 through 7 days. At 6 consecutive days (the legal maximum), the validator permits the assignment. At 7 days, it rejects it. The recursive streak-counting implementation is verified for correctness across non-contiguous assignment histories.

\subsection{Scoring Function}

The \texttt{computeScore} test module verifies each scoring component independently:
\begin{itemize}
    \item A talent with 8 hours remaining scores higher on \texttt{\_score\_hours\_balance} than one with 2 hours remaining.
    \item A talent with a 5-day streak scores $-10.0$ on \texttt{\_score\_streak}; one with no streak scores $0.0$.
    \item A talent whose previous shift ended 10 hours ago scores $-5.0$ on \texttt{\_score\_rest\_gap}.
\end{itemize}

\subsection{Round-Robin Fairness}

The round-robin picker is tested to confirm that across $n$ identical-score candidates, each candidate is selected exactly once before any is selected twice—verifying proportional fairness.

\section{Test Results}

All unit tests pass successfully. The test suite can be executed with:

\begin{lstlisting}[language=bash, style=pythonstyle, caption={Running the test suite}, label={lst:pytest}]
pytest tests/ -v
\end{lstlisting}

Representative output:
\begin{lstlisting}[language=bash, style=pythonstyle, caption={Pytest output summary}]
tests/unit/test_validators.py::test_rest_validator_no_prior   PASSED
tests/unit/test_validators.py::test_rest_validator_11h_exact  PASSED
tests/unit/test_validators.py::test_max_hours_at_limit        PASSED
tests/unit/test_validators.py::test_consecutive_6_days        PASSED
tests/unit/test_validators.py::test_consecutive_7_days        PASSED
tests/unit/test_calculate_score.py::test_hours_balance        PASSED
tests/unit/test_calculate_score.py::test_streak_5_days        PASSED
tests/unit/test_round_robin.py::test_round_robin_fairness     PASSED
tests/unit/test_understaffed.py::test_partial_staffing        PASSED

========= 9 passed in 0.42s =========
\end{lstlisting}

\section{Performance Considerations}

For a typical restaurant scheduling scenario (20 talents, 84 shift slots per week—3 periods × 6 roles × 7 days), the scheduling engine completes in under 200 milliseconds on standard server hardware. This is well within the interactive response time threshold of 1 second \cite{nielsen_1993} for a web API, ensuring that managers experience near-instantaneous schedule generation.

Memory usage is dominated by the \texttt{assignments} list, which grows linearly with the number of assignments and is bounded by the product of talents and days per week. For realistic deployment sizes, peak memory consumption remains well below 50 MB.

% =============================================================
%           CHAPTER 7 — CONCLUSION AND FUTURE WORK
% =============================================================
\chapter{Conclusion and Future Work}

\section{Summary of Contributions}

This thesis has presented SlotMeIn, a production-ready AI-ready workforce scheduling system for hospitality environments. The key contributions are:

\begin{enumerate}
    \item \textbf{A comprehensive literature review} covering workforce scheduling optimisation, constraint satisfaction, fairness in scheduling, predictive analytics for demand forecasting, and AI applications in the hospitality industry. The review identifies a gap in AI-assisted operational scheduling and positions SlotMeIn as a contribution to addressing that gap.

    \item \textbf{A modular RESTful API architecture} developed with FastAPI, SQLAlchemy, and PostgreSQL, following Domain-Driven Design principles. The architecture cleanly separates concerns across bounded contexts and is designed from the ground up for the integration of AI components without architectural disruption.

    \item \textbf{A deterministic constraint-satisfaction scheduling engine} that produces legally compliant, fairness-aware weekly shift plans. The engine implements a least-flexible-first constructive heuristic, a multi-dimensional scoring function (hours balance, consecutive streak, rest gap), and a round-robin tiebreaker—producing consistent, explainable, and auditable scheduling decisions.

    \item \textbf{A validated unit test suite} that systematically verifies the correctness of all four hard constraint validators, the scoring function, the round-robin picker, and the understaffing detector across normal, boundary, and violation scenarios.

    \item \textbf{A data simulation and AI readiness framework} that constructs a labelled dataset from real hotel booking data, performs exploratory data analysis identifying key demand drivers, and defines a machine learning pipeline for per-shift staffing prediction using Gradient Boosting models.
\end{enumerate}

\section{Research Objectives Revisited}

The thesis has fully addressed RO1 (literature review), RO2 (scheduling engine), RO3 (API architecture), and RO4 (unit testing). RO5 (AI predictive layer) has been partially addressed: the data pipeline, EDA, and integration architecture are complete, and model training and evaluation are planned for the next phase.

\section{Limitations}

Several limitations of the current system warrant honest acknowledgement:

\begin{itemize}
    \item The staffing prediction labels in the training dataset are simulated rather than derived from real operational records. While the simulation methodology is grounded in the domain literature, real-world model performance may differ once actual staffing data is available.
    \item The deterministic allocation heuristic is greedy and non-backtracking, which means it cannot guarantee globally optimal solutions. For very constrained instances, this may result in preventable understaffing.
    \item The system currently has no native support for employee self-service—shift preferences, swap requests, and availability updates are managed by the manager rather than directly by employees. This is planned as a future feature.
    \item The frontend client and its UX design are outside the scope of this thesis, though they form an integral part of the deployed SlotMeIn application.
\end{itemize}

\section{Future Work}

The roadmap for continued development of SlotMeIn encompasses the following directions:

\paragraph{AI Predictive Staffing Model.} The primary next step is the training, evaluation, and integration of the Gradient Boosting staffing prediction model described in Chapter~5. Once trained, the model will be evaluated against the deterministic baseline using KPIs including prediction accuracy (RMSE, MAE), reduction in understaffed incidents, and reduction in over-staffed shift hours.

\paragraph{Reinforcement Learning for Schedule Optimisation.} The comment in the allocation engine source code—``state space representation, DFS, BFS, graph and tree search algorithms, neural networks, reinforcement learning''—reflects the author's awareness of the long-term potential for applying reinforcement learning to the scheduling problem. An RL agent that learns scheduling policies from operational feedback (service quality ratings, employee satisfaction scores) represents a compelling long-term research direction \cite{sutton_barto_2018}.

\paragraph{Employee Self-Service.} Extending the system to allow employees to submit availability updates, request time off, and propose shift swaps would significantly reduce manager administrative load and improve employee engagement.

\paragraph{Multi-Venue Support.} The current system models a single venue. Extending it to support multi-location hotel groups—with shared talent pools and cross-venue availability—would substantially broaden its applicability.

\paragraph{Real-Time Demand Integration.} Integrating with live booking management systems (e.g., Opera PMS, Lightspeed) would allow the staffing predictor to consume real-time demand signals rather than relying on historical pattern extrapolation.

\paragraph{Explainability and Transparency.} As the system transitions toward AI-driven staffing decisions, ensuring that managers can understand and interrogate the model's predictions becomes critical. SHAP (SHapley Additive exPlanations) \cite{lundberg_lee_2017} values will be computed and exposed through the API to support explainable AI governance.

\section{Final Remarks}

The development of SlotMeIn has demonstrated that a well-architected, constraint-aware scheduling system—built on established software engineering principles and informed by the operations research literature—can deliver significant value to hospitality organisations. The system is not merely a research prototype; it is a deployed, live application that can be used by real managers today. At the same time, it is designed to grow: the modular architecture, the comprehensive test suite, and the AI readiness framework collectively ensure that SlotMeIn is well-positioned to evolve into the fully AI-powered scheduling intelligence platform envisioned in this thesis plan.

The most important lesson from this project is that intelligence in a scheduling system does not require sacrificing transparency or compliance. By separating the hard constraint layer from the soft scoring layer, and by making the algorithmic logic explicit and testable, SlotMeIn achieves a level of reliability and auditability that pure ML-black-box systems cannot match—while laying the groundwork for those systems to enhance, rather than replace, the deterministic foundation.

% =============================================================
%              BIBLIOGRAPHY
% =============================================================
\printbibliography[heading=bibintoc, title={Bibliography}]

% =============================================================
%              APPENDICES
% =============================================================
\appendix

\chapter{Key Source Code Listings}

\section{ScheduleBuilder — generate\_schedule()}

The following listing presents the complete \texttt{generate\_schedule()} method of the \texttt{ScheduleBuilder} class, which orchestrates the entire allocation pipeline.

\begin{lstlisting}[caption={ScheduleBuilder.generate\_schedule() — full implementation}, label={lst:full_engine}]
def generate_schedule(self):
    plan = []
    working_assignments = list(self.history)

    availability_service = TalentAvailabilityService(
        self.availability, self.assignable_shifts,
        self.talents_to_assign
    )
    eligibility = availability_service.generate_eligible_talents()

    sorted_shifts = sorted(
        self.assignable_shifts.items(),
        key=lambda shift: len(eligibility.get(shift[0], []))
    )

    validators = [
        dailyAssignmentValidator(), restValidator(),
        maxHoursValidator(), consecutiveValidator()
    ]
    round_robin = roundRobinPicker()
    workload = {tid: 0.0 for tid in self.availability.keys()}

    for shift_instance_id, shift in sorted_shifts:
        candidates = eligibility.get(shift_instance_id, [])
        num_assigned = 0
        scorer = computeScore(
            shift=shift,
            availability=self.availability,
            assignments=working_assignments,
            workload=workload
        )
        scores = {tid: scorer.calculate_score(tid)
                  for tid in candidates}
        remaining = set(candidates)

        while num_assigned < shift.role_count and remaining:
            top_score = max(
                score for tid, score in scores.items()
                if tid in remaining
            )
            top_candidates = [
                tid for tid in remaining
                if scores[tid] == top_score
            ]
            best_fit = round_robin.pickBestFit(
                shift.role_name, top_candidates
            )
            if best_fit is None:
                break
            remaining.discard(best_fit)
            ctx = context.contextFinder(
                best_fit, shift,
                self.availability, working_assignments
            )
            if all(v.can_assign_shift(ctx) for v in validators):
                new_assignment = assignment(
                    talent_id=best_fit,
                    shift_id=shift_instance_id,
                    shift=shift
                )
                plan.append(new_assignment)
                working_assignments.append(new_assignment)
                shift_hours = (
                    (shift.end_time - shift.start_time)
                    .total_seconds() / 3600
                    - get_break_duration(shift.shift_name)
                )
                workload[best_fit] += shift_hours
                for v in validators:
                    if hasattr(v, "mark_assigned"):
                        v.mark_assigned(ctx)
                num_assigned += 1

    return plan
\end{lstlisting}

\chapter{Database Entity Relationship Overview}

The following table summarises the foreign key relationships between the seven database tables in SlotMeIn.

\begin{longtable}{p{4cm} p{3.5cm} p{5cm}}
\caption{Database Foreign Key Relationships} \label{tab:fk} \\
\toprule
\textbf{Table} & \textbf{References} & \textbf{Relationship} \\
\midrule
\endfirsthead
\multicolumn{3}{c}{\tablename\ \thetable{} -- continued} \\
\toprule
\textbf{Table} & \textbf{References} & \textbf{Relationship} \\
\midrule
\endhead
\texttt{talent\_constraints} & \texttt{talents} & Many-to-one (one talent, many constraints) \\
\texttt{constraint\_rules} & \texttt{talent\_constraints} & Many-to-one \\
\texttt{shift\_templates} & \texttt{shift\_periods} & Many-to-one \\
\texttt{scheduled\_shifts} & \texttt{talents} & Many-to-one \\
\texttt{scheduled\_shifts} & \texttt{schedules} & Many-to-one \\
\bottomrule
\end{longtable}

\end{document}
